\section{A lower bound in every fixed window}\label{sec:windows}
The bounded one-bit response $4K$ alone cannot force visits to windows narrower than that bound. A long plus buffer reduces the response of its centre to any prescribed positive size. The derivative remains macroscopic because a positive density of centres have such buffers.

\begin{lemma}[Arbitrarily small buffered changes]\label{lem:buffer}
For every $\eps>0$ there is $L<\infty$ such that, for all sufficiently large $n$, a centre in the middle half of the north collar with $L$ plus bits on each side satisfies
\[
 0\le\osc_iF_n<\eps
\]
uniformly over all the other exterior spins.
\end{lemma}
\begin{proof}
Compute the response with the centre itself plus. Directly inside its $2L+1$ plus wall bits take the half-box
\[
 V_L=\{(x+j,n-k): -L\le j\le L,\ 0\le k\le L\}.
\]
For fixed $L$ and large $n$ this box avoids $J_n$. Condition either pinned ensemble on spins outside $V_L$. Ferromagnetic monotonicity brackets the conditional law by the two laws on $V_L$ with the same plus north wall and with the other three exterior sides all plus or all minus. Let $e_L$ be the difference of their plus-probabilities at $y=(x,n)$.

As $L\to\infty$, any subsequential limit of either extremal half-box law satisfies the half-plane specification with the uniform plus wall. For every finite set this follows by passing its finite-range DLR identity to the limit; the other three sides recede to infinity. Uniform-wall uniqueness, with $h=1$ as checked in Section~\ref{sec:model}, makes all these limits identical. Compactness therefore implies $e_L\to0$. The bracketing is independent of the conditioned surroundings, pins, and distant exterior bits, proving the required uniformity.

With the centre plus, changing it to minus tilts by $e^{-2K\sigma_y}$. The absolute derivative of
\[
 \log[e^{2K}+(e^{-2K}-e^{2K})u]
\]
in the one-site plus-probability $u$ is at most $e^{4K}-1$. The two pinned marginals at $y$ differ by at most $e_L$, so $\osc_iF_n\le(e^{4K}-1)e_L$. Choose $L$ with the last quantity strictly below $\eps$.
\end{proof}

\begin{lemma}[Drift and variance for buffered centres]\label{lem:blockdrift}
Fix $L$ and use the block law $\Pp_q^{(L)}$ of Section~\ref{sec:contours}. There are $c_L>0$ and $C_L<\infty$ such that, uniformly for $|q-1/2|\le n^{-2/5}$ and all large $n$,
\begin{equation}\label{eq:blockdrift}
 \frac d{dq}\E_q^{(L)}F_n\ge c_Ln,\qquad
 \Var_q^{(L)}(F_n)\le C_Ln.
\end{equation}
Also $\E_{1/2}^{(L)}F_n=0$.
\end{lemma}
\begin{proof}
There are $N_0=\lfloor(2\lfloor n/2\rfloor+1)/b\rfloor$ complete blocks, where $b=2L+1$, so $N_0\ge n/(2b)$ eventually. A given centre is active with probability $a=2^{-2L}$, independently of $q$. The concentration and surface-separation estimates were proved for these independent blocks, including partially intersected blocks, in Lemmas~\ref{lem:ratios} and \ref{lem:tooth}.

Choose the depth $D=D(L)$ so large that each of the two expected remainders at every centre is eventually below $ag/8$, uniformly over the specified $q$-window. From \eqref{eq:gap} and nonnegativity of the remainders,
\begin{equation}\label{eq:activegap}
 \E_q^{(L)}[\one_{\{i\text{ active}\}}d_v]
 \ge ag-\E_q^{(L)}R_{n,v}^+-\E_q^{(L)}R_{n,v}^-
 \ge\frac{3ag}{4}.
\end{equation}
This estimate is unconditional; activity is not imposed on a contour estimate. By \eqref{eq:response}, the mean active response is at least $3ag\gamma_D/4$. Conditional on buffers and inactive centres, differentiation in $q$ sums the oscillations over the active centres, whose bits are independent Bernoulli($q$). Averaging gives
\[
 \frac d{dq}\E_q^{(L)}F_n
 \ge N_0\frac{3ag\gamma_D}{4}
 \ge\frac{3ag\gamma_D}{8b}n.
\]
Thus $c_L=3ag\gamma_D/(8b)$ is permissible.

Expose the independent whole blocks and singletons. Changing one group changes $F_n$ by at most $4Kb$ and there are at most $9n$ groups. The variance bound in \eqref{eq:bd} gives $C_L=36K^2b^2$. At $q=1/2$ the whole word is fair, so oddness gives mean zero. In particular neither $c_L$ nor $C_L$ depends on a subsequent choice of a tilt constant.
\end{proof}

\begin{proposition}[A window for the pinned odds]\label{prop:oddswindow}
For every $z\in\R$ and $\eps>0$, there are $c(z,\eps)>0$ and $n_0$ such that
\begin{equation}\label{eq:oddswindow}
 \Pp_{1/2}(|F_n-z|<\eps)\ge c(z,\eps)n^{-1/2}
 \qquad(n\ge n_0).
\end{equation}
\end{proposition}
\begin{proof}
Choose $L$ by Lemma~\ref{lem:buffer} for jumps strictly below $\eps$, and fix the constants in Lemma~\ref{lem:blockdrift}. Take
\[
 A=1+\frac{16\sqrt{C_L}}{c_L},\qquad
 q_\pm=\frac12\pm\frac A{\sqrt n}.
\]
For large $n$ these lie in $(0,1)$ and satisfy $|q_\pm-1/2|\le n^{-2/5}$. Integrating \eqref{eq:blockdrift} from $1/2$ puts the endpoint means respectively above $c_LA\sqrt n$ and below $-c_LA\sqrt n$. Once $|z|\le c_LA\sqrt n/2$, Chebyshev gives
\begin{equation}\label{eq:straddle}
 \Pp_{q_+}^{(L)}(F_n\le z)<1/16,\qquad
 \Pp_{q_-}^{(L)}(F_n\ge z)<1/16,
\end{equation}
since $4C_L/(c_L^2A^2)<1/64$.

Couple the whole family of laws by fixing all buffers and fair inactive bits, and taking independent uniform variables $U_i$ at the active centres: the centre is plus exactly when $U_i\le q$. Conditional on the buffers, let $N$ be their number and $S(q)$ their plus count. At each endpoint $S(q_\pm)$ has law $\operatorname{Bin}(N,q_\pm)$, where $N\le9n$. Its mean differs from $N/2$ by at most $9A\sqrt n$ and its variance is at most $9n/4$. With $B_*=9A+12$, Chebyshev and a union bound give
\begin{equation}\label{eq:countband}
 \Pp\left(\max_{q\in\{q_-,q_+\}}|S(q)-N/2|>B_*\sqrt n
       \,\middle|\,\text{buffers}\right)\le1/32.
\end{equation}
Consequently with probability at least $1/2$ the endpoint values strictly straddle $z$ and both counts lie in this band. All intermediate counts lie in it too, by monotonicity.

Almost surely the finitely many thresholds are distinct. On this event take the first flip in $(q_-,q_+)$ taking $F_n$ from a value at most $z$ to a value above $z$. It is an active flip, so Lemma~\ref{lem:buffer} makes both adjacent values lie strictly in $(z-\eps,z+\eps)$. This remains true if the preceding value equals $z$. Record such crossing flips only when their pre-flip count lies in \eqref{eq:countband}. Let $V$ be their number. The preceding event implies $\E V\ge1/2$.

We bound $\E V$ by integration of a single uniform threshold. Fix a potential centre $i$ and its buffers. If it is active and $U_i=q$, all the other active centres at that threshold are independent Bernoulli($q$). If their plus count is $s$, the pre-flip count condition implies
\[
 |s-(N-1)/2|\le B_*\sqrt n+1.
\]
Their probability density relative to independent fair signs is
\[
 R_q=(2q)^s[2(1-q)]^{N-1-s}.
\]
Writing $\delta=q-1/2$ and using $\log(1+x)\le x$ gives
\begin{equation}\label{eq:likelihood}
 \log R_q\le 4\delta[s-(N-1)/2]\le4A(B_*+1)
\end{equation}
throughout the recorded band. This bound is uniform in the buffer word and in $N$, including small $N$. All inactive bits and buffers have the same distribution for every $q$.

The crossing criterion depends only on the other bits, with bit $i$ assigned minus before and plus after the flip. On that criterion \emph{both} assignments have $|F_n-z|<\eps$. Under fair sampling of the other bits its probability, including the activity requirement, is therefore at most the unconditional fair window probability: averaging also over the independent fair bit $i$ preserves the window. There is no conditioning cost for any particular buffer pattern. Bound the tilted criterion by \eqref{eq:likelihood}, integrate $q$ over its interval of length $2A/\sqrt n$, and sum over at most $9n$ potential centres. It follows that
\[
 \frac12\le\E V
 \le18A e^{4A(B_*+1)}\sqrt n\,
       \Pp_{1/2}(|F_n-z|<\eps).
\]
Thus $[36A e^{4A(B_*+1)}]^{-1}$ is a permissible positive constant in \eqref{eq:oddswindow} once the large-size cutoffs above hold. Its size is immaterial for recurrence.
\end{proof}

\begin{proof}[Proof of Theorem~\ref{thm:lower}]
Choose $z\in\R$ and $\eps>0$ so that the compact interval
$\logistic([z-\eps,z+\eps])$ lies strictly inside $(u,v)$. On $H_n$, \eqref{eq:transfer} places $T_n$ in $(u,v)$ whenever $|F_n-z|<\eps$, for all sufficiently large $n$. The event $H_n^c$ has probability $o(n^{-1/2})$ by \eqref{eq:H}, so Proposition~\ref{prop:oddswindow} proves the lower bound after reducing its constant.

If $U$ is a weakly open neighbourhood of $\nu_\lambda$ with $0<\lambda<1$, continuity of $\theta\mapsto\nu_\theta$ provides $0<u<v<1$ and a positive distance margin such that $\lambda\in(u,v)$ and $\nu_\theta\in U$ with that margin for $\theta\in[u,v]$. Lemma~\ref{lem:calibration} then puts $\mu_n^\eta$ in $U$ whenever $u<T_n<v$ and $n$ is large. Apply the already proved bound.
\end{proof}
